% Cuatro niveles de diagnóstico del daño (Rytter): cada nivel incluye al anterior.
\documentclass[tikz,border=4pt]{standalone}
\usepackage{resonancia}
\begin{document}
\begin{tikzpicture}
  \foreach \i/\t/\q/\c in {1/Detección/{¿Hay daño?}/rAzul!12, 2/Localización/{¿Dónde está?}/rAzul!22, 3/Cuantificación/{¿Qué gravedad tiene?}/rAzul!34, 4/Pronóstico/{¿Cuánto servicio queda?}/rAzul!48} {
    \pgfmathsetmacro{\x}{(\i-1)*3.55}
    \pgfmathsetmacro{\h}{0.9+0.75*\i}
    \draw[trazo, color=rAzul, fill=\c, rounded corners=2pt] (\x,0) rectangle (\x+3.45,\h);
    \node[font=\sffamily\bfseries\small, color=rTinta, anchor=north west] at (\x+0.12,\h-0.1) {\i\ · \t};
    \node[etiqueta, anchor=south west, color=rTinta, font=\sffamily\footnotesize] at (\x+0.12,0.12) {\q};
  }
  \draw[flecha, color=rNaranja] (0.1,-0.35) -- (13.9,-0.35) node[midway, below, etiqueta, color=rNaranja] {cada nivel incluye al anterior y exige más información y mejores modelos};
\end{tikzpicture}
\end{document}
